\section{Power on clean data}\label{sec:power}

\subsection{Level and choice of basis}\label{sec:size}\label{sec:basis}

Table~\ref{tab:size} gives the size of each basis and weighting over the $131$ null scenarios. The
median is $0.048$ to $0.050$ in every form and the range is $0.013$ to $0.064$. In the unit form the
cubic basis falls outside a band of three Monte Carlo standard errors in one scenario of $131$, at the
$5\%$ level, the symmetric basis in one, at the $1\%$ level, and the Stukel basis in three, all at
$1\%$. The score form is less well behaved, failing four scenarios with the cubic basis and ten with
the Stukel basis. The rule fixed in advance therefore selected the unit form for all three bases.

\input{tab_size}

Table~\ref{tab:families} gives the mean size-adjusted power of the three bases in each departure
family. No basis leads everywhere. The cubic basis leads on asymmetric links, by $0.25$ over the
symmetric basis, and on the off-index scenarios; the symmetric basis leads on symmetric tails and on
rough misfit. Averaged over the six families the cubic basis leads narrowly, $0.452$ against $0.442$
and $0.425$, so it is the default, and an analyst who knows which departure they fear should choose
the basis from the row for that family rather than from the average.

\input{tab_families}

\subsection{The partition, Stukel's test and the calibration belt}\label{sec:verdicts}

\paragraph{Refining the partition} On clean data the default partition gives EDGE more
power than ten groups: it leads or ties the ten-group form in $155$ of the $158$ scenarios, with a mean
gain of $0.017$ in size-adjusted power. Against the partition tests the gain is larger, because the
Hosmer--Lemeshow statistic cannot use a finer partition without spending a degree of freedom on each
new group. Against the best of four partition tests (the Hosmer--Lemeshow statistic, its Farrington
correction, its equal-width form and the Pigeon--Heyse test), EDGE at the default
partition leads or ties in all $21$ scenarios that the specification selected as detectable and not
saturated; with ten groups for every test, in $17$ of the $19$ selected for that comparison.

\paragraph{Against Stukel's test} Stukel's score test is the established directed test for this
problem. The default basis stays within $0.02$ of Stukel's joint score in mean size-adjusted power in
every departure family, and over all $158$ scenarios it is $0.011$ behind (Monte Carlo $95\%$ interval
$0.010$ to $0.013$). On
the symmetric-tail family, the departure Stukel's one-parameter score was built for, the one-parameter
score is ahead of the default basis by $0.062$; there the symmetric basis is the one to use,
and in its score form it matches the one-parameter score, $0.002$ behind on average. In its unit form it
is level with Stukel's joint score on the same family, $0.005$ ahead. On clean data EDGE is
therefore comparable to Stukel's test; they differ in how they behave under corrupted records
(Section~\ref{sec:contamination}).

\paragraph{Against the calibration belt} On the symmetric-tail family the cubic basis leads the GiViTI
belt by $0.105$ on average, but not in every regime: in the scenarios with $12\%$ events the belt leads,
by $0.152$ and $0.124$ in the two probit scenarios and by smaller margins in four others. On the
asymmetric-link family the cubic basis is $0.075$ behind the belt. Both results have the same cause.
When events are scarce, or when the calibration curve bends once and smoothly, a polynomial fitted to
the whole risk scale uses information that a partition averages away. These are regimes in which the
belt, or for the single bend the cubic calibration test, is the better test, and
Section~\ref{sec:whennot} lists them as such.

\subsection{Against the smoothing and resampling tests}\label{sec:rivals}

Table~\ref{tab:rivals} gives raw rejection rates of le Cessie's smoothed test, the Liu projection test
and BAGofT beside EDGE on the $20$ link and tail scenarios at $n\ge500$. Le Cessie's test
holds its level in all ten matched null scenarios, between $0.040$ and $0.064$. On size-adjusted power
over the same $500$ replicates, EDGE leads or ties le Cessie's test in $19$ of the $20$
scenarios, is ahead after the Holm correction in $14$, and has mean power $0.451$ against $0.314$. It
leads the projection test in $17$ scenarios and BAGofT in $18$, by up to $0.364$ and $0.630$. All three
rivals lead it clearly in one scenario, the same one for all three: the crossover departure, in which
the calibration curve crosses the diagonal twice. There le Cessie's test rejects $496$ of $500$
replicates and the projection test $478$, against EDGE's $37$, and BAGofT rejects $155$ of
$200$ against $18$.

\input{tab_rivals}

The same three tests were run on the families they are built for, omitted terms and rough misfit:
le Cessie's test on the $56$ such scenarios with $n\le1000$, the projection test on the $19$ at
$n=500$ and BAGofT on six of those. Every matched null held its level, between $0.025$ and $0.070$.
Where the omitted term acts through the fitted risk, a quadratic term, EDGE leads or ties
le Cessie's test in $18$ of $21$ scenarios; where it does not, an interaction, le Cessie's test leads in
$19$ of $27$, and the projection test on all four continuous interactions. On the fastest
oscillation and the sawtooth BAGofT reaches $1.000$ and $0.805$ against EDGE's $0.295$ and $0.045$. A
basis of smooth shapes cannot follow a fast oscillation, and a test that learns its own partition can.

At $n=200$, where the default partition is ten groups and there is nothing to refine,
EDGE leads or ties le Cessie's test and the projection test in all six link and tail scenarios and
BAGofT in five. Its mean size-adjusted power is $0.160$ against $0.111$ and $0.125$, and $0.098$
against BAGofT's $0.073$ on the replicates BAGofT was run on. Two paired differences survive the Holm
correction, both in its favour and both on the asymmetric Stukel design, where its size-adjusted power
is $0.450$ against le Cessie's $0.330$ and $0.350$ against BAGofT's $0.185$.

\subsection{All scenarios, and what each test costs}\label{sec:census}\label{sec:compute}

Table~S7 of the Supporting Information gives, for each departure family, the mean difference in
size-adjusted power between EDGE at its default setting and each comparator, with a Monte Carlo $95\%$
interval. Over all $158$ scenarios EDGE is ahead of the Hosmer--Lemeshow statistic at ten groups by
$0.064$ ($0.063$ to $0.066$): by $0.18$ on asymmetric links and symmetric tails, $0.13$ on the Stukel
symmetric family and $0.07$ on omitted terms, and behind by $0.46$ on rough misfit, which a basis of
smooth shapes cannot follow. It is level with Stukel's joint score ($-0.011$), the one-parameter score
($-0.007$) and the cubic calibration test ($-0.008$), and $0.015$ ahead of the GiViTI belt, which leads by
$0.07$ on asymmetric links and trails by $0.10$ on symmetric tails. The tie with the cubic calibration
test is expected, because EDGE's default basis looks for the same shapes: a cubic in the predicted risk,
where that test fits a cubic in the logit. The two tests ask the same question of clean data; they
differ in how a record enters the answer, and that is what decides their behaviour when a few records
are wrong (Section~\ref{sec:whatprotects}). Figure~\ref{fig:scenarios} shows the
same comparison scenario by scenario and Figure~\ref{fig:census} by family; the counts in their strips
use a margin of $0.01$, within which two tests cannot be told apart. Against the other grouped tests
EDGE is behind in $14$ of the $80$ omitted-term scenarios, $10$ of the $13$ rough-misfit scenarios and
both off-index scenarios; $20$ of these $26$ losses are to Xie's and Tsiatis's tests, which group on the
covariates.

\begin{figure}[htbp]\centering
\includegraphics[width=\textwidth]{Fig13_scenarios}
\caption{EDGE at its default setting against every comparator, one point per alternative scenario. Both
axes are size-adjusted power, so each test is read at its own realised level; a point in the green
half is a scenario in which EDGE is ahead. The strip gives the number of scenarios in
which it leads or ties, within $0.01$. The first panel takes, in each scenario, the better of the
Hosmer--Lemeshow statistic and its Farrington correction at ten groups; the second, the best of the
equal-width Hosmer--Lemeshow statistic, the Pigeon--Heyse, Tsiatis and Xie tests. Le Cessie's test and
the two resampling tests were run on a subset of the scenarios, and the strip gives how many. Colour
is the departure family.}
\label{fig:scenarios}
\end{figure}

\begin{figure}[htbp]\centering
\includegraphics[width=\textwidth]{Fig14_census}
\caption{The comparison of Figure~\ref{fig:scenarios} by departure family: each entry is the number of
scenarios in which EDGE leads or ties the comparator in that row. The last column is the
total over all families.}
\label{fig:census}
\end{figure}

Figure~\ref{fig:cost} sets power on clean data against protection and cost, on one core with nothing
else running. At $n=1000$ EDGE takes $4$ milliseconds, Stukel's joint score $1$, the
Hosmer--Lemeshow statistic $2$, the cubic calibration test $5$ and the GiViTI belt $9$; le Cessie's test
takes $5.6$ seconds, the projection test $35$ and BAGofT $392$, three to five orders of magnitude more.
On the twenty scenarios of panel (a) the cubic calibration test ($0.477$) is slightly ahead of EDGE
($0.440$) and Stukel's joint score ($0.446$) level with it, but neither survives one corrupted record in a thousand
(Section~\ref{sec:contamination}). Among the protected tests EDGE has the most power: at ten groups
$0.411$, against $0.267$ for the Hosmer--Lemeshow statistic, the only cheaper one, and $0.314$ for
le Cessie's test, at more than a thousand times the cost in its published form. These panel (a) values
use each test's full replicates; on the $500$ replicates of Table~\ref{tab:rivals} EDGE reads $0.451$
(Section~\ref{sec:rivals}).

The tests without resampling grow roughly in proportion to $n$: at $n=100{,}000$ EDGE takes half
a second, more than Stukel's score and the Hosmer--Lemeshow statistic ($18$ and $120$ milliseconds)
because it also builds the covariance of its partition. Le Cessie's test in its published form multiplies $n\times n$
matrices and grows with the cube of $n$: $74$ seconds at $n=2000$ and $21$ minutes at $n=5000$, where EDGE
takes $14$ milliseconds. The rates reported for le Cessie's test were computed with its implementation
in \texttt{ebrahim.gof}, which gives the same statistic. BAGofT's cost is fixed by its design rather than by $n$: $(1+\textit{nsim})\times
\textit{nsplits}$ model fits, ten minutes for one test at $n=2000$.

\begin{figure}[htbp]\centering
\includegraphics[width=\textwidth]{Fig15_cost}
\caption{Power, protection and cost of each test. (a)~Mean size-adjusted power on clean data
over the twenty scenarios on which all eight tests were run (link and tail departures and one
crossover), against the false-alarm rate when ten records in $1000$ carry a corrupted covariate and the
model is otherwise correct (corruption C1, Section~\ref{sec:contamination}); the shaded strip is below
$0.10$. Each label gives the median time for one data set at $n=1000$. The two resampling tests were
not run with corrupted records. The power axis ranks the tests on one family only, the one in which a
curve fitted to the whole risk scale is strongest; over all $158$ scenarios EDGE leads or
ties the cubic calibration test in $123$ (Figure~\ref{fig:scenarios}), as expected of two tests that
look for the same cubic shapes (Section~\ref{sec:census}). (b)~Median time against sample size, both axes
logarithmic, on one core. Each rival is timed as published: BAGofT from its own package at its
defaults, whose cost is $10{,}100$ model fits, le Cessie's test in its published $O(n^3)$ form, and the
projection test from our implementation at $B=250$. The tests without resampling are timed over eleven
repetitions, le Cessie's test over three at $n\le1000$; le Cessie's test at $n\ge2000$, the projection
test at $n=1000$ and BAGofT are timed once.
Colour gives the cost class in (a); EDGE is shown at the default partition and at ten groups.}
\label{fig:cost}
\end{figure}
